\chapter{Testing}
\label{ch:testing}

\section{Current state}

\textbf{Status:} NOT IMPLEMENTED (no automated tests exist).

A recursive search of the working tree for test or specification files
(\texttt{*.test.ts}, \texttt{*.spec.ts} and JavaScript equivalents) returns no
results, and no test runner is declared in any package. The root
\texttt{package.json} defines a \texttt{test} script (\texttt{pnpm -r test}),
but no workspace package defines a \texttt{test} script for it to run. The VS
Code extension README lists ``Testing infrastructure'' as an unfinished roadmap
item and ``No tests'' under known limitations, and the CHANGELOG contains no
test-related entries.

What \emph{does} exist and acts as a lightweight correctness gate:

\begin{itemize}
  \item \texttt{scripts/validate.mjs} --- a recipe validator (see below).
  \item \texttt{pnpm build} across all packages, which performs TypeScript
        compilation under \texttt{strict} mode and therefore catches type
        errors.
  \item \texttt{pnpm --filter expo-snips typecheck}, an explicit no-emit type
        check for the extension.
  \item \texttt{scripts/build-registry.mjs}, which fails if a recipe references a
        missing template.
\end{itemize}

These are build-time checks, not tests. They do not execute behaviour.

\section{The validator as the de-facto quality gate}

\texttt{scripts/validate.mjs} is the closest thing to an automated test. It is
an inline, hand-written validator (it reproduces a subset of Draft-07 JSON
Schema rather than using a schema library) and checks each recipe for:

\begin{longtable}{@{}L{4.6cm}L{9.6cm}@{}}
\caption{Checks performed by \texttt{validate.mjs}}\\
\toprule
\textbf{Check} & \textbf{Detail} \\
\midrule
\endfirsthead
\toprule
\textbf{Check} & \textbf{Detail} \\
\midrule
\endhead
Required fields & Every field in \texttt{schema.required} must be present and
non-null. \\
Id format & Must match \texttt{\^{}[a-z0-9-]+/[a-z0-9-]+\$}. \\
Version format & Must match \texttt{\^{}\textbackslash d+\textbackslash.\textbackslash
d+\textbackslash.\textbackslash d+\$}. \\
Enums & \texttt{framework} must be \texttt{react-native}; \texttt{runtime} must be
\texttt{expo}. \\
Supported option values & Each entry of \texttt{supported.styling/router/ui/icons}
must be one of the schema enums. \\
Files present & \texttt{files} must be a non-empty array; each entry needs
\texttt{path} and \texttt{templateFile}. \\
Id matches path & The recipe id must equal its path relative to
\texttt{data/} (cross-check between content and location). \\
Template existence & Every referenced \texttt{.hbs} file must exist on disk. \\
\bottomrule
\end{longtable}

\textbf{Observations:}
\begin{itemize}
  \item The validator is genuinely useful: it enforces identity, structure and
        referential integrity of the catalog.
  \item It does \emph{not} validate the \texttt{variants}, \texttt{primitives},
        \texttt{tags} or \texttt{dependencies} fields beyond their presence, and
        it does not validate nested \texttt{when} conditions.
  \item Because it re-implements schema validation by hand, it can drift from
        \texttt{schema.json}; the schema itself is not executed.
  \item It is invoked by the developer (or by the bake pipeline), not wired into
        CI (see Chapter~\ref{ch:cicd}).
\end{itemize}

\section{Rendering is untested}

The generation engine (\texttt{packages/core}) --- the part most likely to
regress and the part both clients depend on --- has no tests. There is no
assertion that a given recipe, options set and variant produce the expected
file paths and content, no test of the \texttt{when} filtering logic, and no
test of primitive recursion. A change to a template or to
\texttt{getFilesForRecipe} is verified only by building and, at best, running the
tool by hand.

\section{Test coverage by area}

\begin{longtable}{@{}L{4cm}L{2.6cm}L{7.6cm}@{}}
\caption{Coverage assessment by area}\\
\toprule
\textbf{Area} & \textbf{Coverage} & \textbf{Comment} \\
\midrule
\endfirsthead
\toprule
\textbf{Area} & \textbf{Coverage} & \textbf{Comment} \\
\midrule
\endhead
Recipe data & Structural only & \texttt{validate.mjs} checks shape and template
existence, not rendered output. \\
Core engine & None & No unit tests for rendering, filtering or recursion. \\
CLI commands & None & No command-level tests. \\
CLI utilities & None & Diff, detection, manifest and install logic untested. \\
VS Code extension & None & No integration or unit tests; no
\texttt{@vscode/test-electron} harness. \\
Scripts & None & The build, bake and serve scripts are not tested. \\
Type safety & Partial & \texttt{strict} TypeScript catches type-level errors at
build. \\
\bottomrule
\end{longtable}

\section{Recommended test strategy}

\textbf{Status:} RECOMMENDED. In order of value per unit of effort:

\begin{enumerate}
  \item \textbf{Golden-file tests for the core engine.} For a representative set
        of recipes and option combinations, assert that \texttt{getFilesForRecipe}
        returns a fixed set of paths and that rendered content matches a stored
        snapshot. This is the highest-value test because it protects the shared
        engine and the templates simultaneously.
  \item \textbf{Variant and \texttt{when} matrix tests.} Assert that each
        supported styling/router/ui/icons combination selects the intended
        template, and that a variant overrides the default.
  \item \textbf{Primitive recursion tests.} Assert that a shadcn recipe pulls in
        the expected primitive files, that existing files are skipped, and that a
        synthetic cycle terminates.
  \item \textbf{Manifest tests.} Assert \texttt{upsertEntry} semantics for the
        same recipe with different variants, and \texttt{updateEntryVersions}
        behaviour.
  \item \textbf{CLI smoke test.} Run \texttt{add} against a temporary directory
        and assert the resulting file tree and manifest, using
        \texttt{--no-install} to avoid network and package-manager side effects.
  \item \textbf{Schema-driven validation.} Replace or supplement the hand-written
        checks with a real JSON Schema validator so \texttt{schema.json} is the
        single source of truth.
  \item \textbf{Extension integration test.} A minimal
        \texttt{@vscode/test-electron} run that activates the extension and
        generates into a scratch workspace.
\end{enumerate}

A suitable stack is already implied by the codebase: \texttt{vitest} or the
Node built-in test runner for the Node packages, and
\texttt{@vscode/test-electron} plus Mocha for the extension (the VS Code
convention). No test framework is currently installed, so this would be a new
dependency.

\section{Summary}

\begin{table}[H]
\centering
\caption{Testing posture}
\begin{tabular}{@{}L{4.4cm}L{9.8cm}@{}}
\toprule
\textbf{Dimension} & \textbf{Assessment} \\
\midrule
Unit tests & None. \\
Integration tests & None. \\
Structural validation & IMPLEMENTED (\texttt{validate.mjs}). \\
Type checking & IMPLEMENTED (\texttt{tsc --strict}, extension \texttt{typecheck}). \\
CI test gate & None (see Chapter~\ref{ch:cicd}). \\
Manual verification & The only current method for behaviour. \\
\bottomrule
\end{tabular}
\end{table}

\textbf{Status:} NOT IMPLEMENTED / RECOMMENDED.
